% Originally in Beginning of Section 5
\textbf{Evaluation Metrics.} We evaluate the correctness of a model response at two levels of granularity. The coarse-grained metric Success Rate, defined as the proportion of task episodes where $reward = 1$, can also be viewed as a stricter form of execution accuracy in that the performed action should not only be valid and yield the intended output but also not result in any unnecessary side-effects. Additionally, we provide more detailed feedback instead of a binary 0/1 assessment. The calculation of rewards for Bash and SQL are extensively discussed in \S\ref{appx:env:bashenv} and \S\ref{appx:env:sqlenv}.
We do not rely on the exact match accuracy of the generated action itself as SQL constructs and bash utilities offer various ways to express the same action, making it challenging for zero-shot models to achieve a high match.

% Originally in Section 5.2: Prompting for Reasoning Comparison
With the ReAct strategy, the agent demonstrates a better comprehension of the intended query and makes accurate decisions on exploring further or executing actions sequentially, particularly in Bash, where it identifies core operations and chains or pipes subsequent operations accordingly. Similarly, when encountering missing bash utilities or \textit{TableName/ColumnName} errors in SQL, the agent chooses to check/install them or acquire more knowledge about the database properties before attempting a retry. On the other hand, the Plan \& Solve strategy formulates a premeditated plan and executes it, although it proves to be less effective in dynamically adapting the plan when errors arise in subsequent steps. These strategies yield success rates of around 60-70\% after interacting for approximately 4-6 turns, highlighting the challenging and open nature of these tasks.